\section{Zadání}
GPS aparaturou změřte zeměpisné souřadnice bodů $\varphi, \lambda$ bodů $P_1, P_2$ vztažené k elipsoidu WGS-84 a vypočtěte jejich obraz v rovině Křovákova zobrazení. Měření proveďte na bodech o známých souřadnicích (trigonometrický bod), a to opakovaně.

Pro transformaci mezi elipsoidy použijte sedmiprvkovou Helmertovu prostorovou podobnostní transformaci danou vztahem

\[
\begin{pmatrix}
X \\
Y \\
Z
\end{pmatrix}
=
m
\begin{pmatrix}
1 & \omega_z & -\omega_y \\
-\omega_z & 1 & \omega_x \\
\omega_y & -\omega_x & 1
\end{pmatrix}
\begin{pmatrix}
X' \\
Y' \\
Z'
\end{pmatrix}
+
\begin{pmatrix}
\Delta X \\
\Delta Y \\
\Delta Z
\end{pmatrix}.
\]

Parametry prostorové transformace:

\[
\begin{array}{|c|c|}
\hline
\omega_x & 4.9984'' \\
\hline
\omega_y & 1.5867'' \\
\hline
\omega_z & 5.2611'' \\
\hline
m - 1 & -3.5623 \times 10^{-6} \\
\hline
\Delta X & -570.8285 \ \mathrm{m} \\
\hline
\Delta Y & -85.6769 \ \mathrm{m} \\
\hline
\Delta Z & -462.8420 \ \mathrm{m} \\
\hline
\end{array}
\]

Určení zeměpisné šířky $\varphi$ obrazu bodu $P$ na Besselově elipsoidu proveďte s přesností $0.001''$. Souřadnice obrazu bodu $P$ v rovině Křovákova zobrazení určete s přesností na cm. V bodech $P_1, P_2$ dále vypočtěte hodnoty délkového zkreslení $m - 1$.

\begin{itemize}
    \item hodnoty délkového zkreslení $m - 1$,
    \item hodnoty meridiánové konvergence $c$.
\end{itemize}

Volitelně porovnejte účinek zanedbání změny/vlivu elipsoidu na souřadnice bodů $P_1, P_2$ a na vzdálenosti $\|P_1 - P_2\|$. Jak se budou lišit takto určené souřadnice a vzdálenosti od hodnot vztažených k elipsoidu WGS-84?

Obrazy bodů $P_1, P_2$ a jejich blízkého okolí vizualizujte s použitím vhodných podkladových geografických dat (např. WMS). Výpočty realizujte v SW Matlab.

\section*{Hodnocení}

\begin{table}[b]
\centering
\begin{tabular}{|l|c|}
\hline
\textbf{Krok} & \textbf{Hodnocení} \\
\hline
Výpočet souřadnic bodu v rovině Křovákova zobrazení + zkreslení. \cellcolor{green} & 20b \cellcolor{green} \\
\hline
\textit{Výpočet meridiánové konvergence} \cellcolor{green} & +5b \cellcolor{green} \\
\hline
\textit{Zanedbání změny elipsoidu: souřadnice $\varphi, \lambda$ měřeny na Besselově elipsoidu,} & +5b \\
\hline
\textit{Zanedbání vlivu elipsoidu: souřadnice měřeny na Gaussově kouli.} & +5b \\
\hline
\textit{Transformace $(\varphi, \lambda, H) \rightarrow (Y, X, Z).$} \cellcolor{green} & +10b \cellcolor{green}\\
\hline
\textit{Použití dotransformace (přesnost 2cm)} & +15b \\
\hline
\textbf{Max celkem:} & \textbf{60b} \\
\hline
\end{tabular}
\end{table}


